\chapter*{Homework 5}
\section*{Problem 1}
Let $\left(U_i\right)_{i \in \mathbb{N}}$ be an i.i.d sequence 
of random variables with $U_i \sim U([0,1])$. Show that
\begin{itemize}
\item $\lim _{n \rightarrow \infty}\left(U_1 U_2 \ldots U_n\right)^{1 / n}=e^{-1}$ almost surely.
\item $\lim _{n \rightarrow \infty} U_1 U_2 \ldots U_n=0$ almost surely.
\end{itemize}
\begin{proof}
\begin{itemize}
\item
Let
\[
X_i := -\log U_i, \quad i \in \mathbb{N}
\]
Since \(U_i \sim U([0,1])\), for \(x \ge 0\),
\[
\mathbb{P}(X_i \le x)
= \mathbb{P}(-\log U_i \le x)
= \mathbb{P}(U_i \ge e^{-x})
= 1 - e^{-x}
\]
Thus \(X_i \sim \mathrm{Exp}(1)\), so
\(
\mathbb{E}[X_i] = 1
\). 
Also,
\[
\log\left((U_1U_2\cdots U_n)^{1/n}\right)
= \frac{1}{n}\sum_{i=1}^n \log U_i
= -\frac{1}{n}\sum_{i=1}^n X_i
\]
By the Strong Law of Large Numbers,
\[
\frac{1}{n}\sum_{i=1}^n X_i = \frac{1}{n}\sum_{i=1}^n \log U_i\to 1
\quad \text{a.s.}
\]
Since the exponential function is continuous,
\[
(U_1U_2\cdots U_n)^{1/n}
= \exp\left(\frac{1}{n}\sum_{i=1}^n \log U_i\right)
\to e^{-1}
\quad \text{a.s.}
\]

\item
Let
\[
P_n := U_1U_2\cdots U_n
\]
Then from the first part we instantly have
\[
P_n^{1/n} \to e^{-1} < 1
\quad \text{a.s.}
\]
Then for any event \(\omega\) in the event of probability one where this convergence holds,
choose \(r\) s.t. 
\(
e^{-1} < r < 1
\), 
then for all sufficiently large \(n\),
\[
P_n(\omega)^{1/n} < r,\quad \text{i.e.} \quad P_n(\omega) < r^n
\]
Since \(0<r<1\), we have \(r^n \to 0\). Therefore
\(
P_n(\omega) \to 0
\). Therefore 
\[
U_1U_2\cdots U_n \to 0
\quad \text{a.s.}
\]
\end{itemize}
\end{proof}



\section*{Problem 2}
A factory produces small resistors, and the resistance of each resistor 
is a random variable $X_i$ with unknown mean $\mu$ and variance $\sigma^2=0.25$ ohms $^2$.
The quality control engineer wants to estimate the average resistance of a batch. 
She decides to measure $n$ resistors and compute the sample average
$$
\bar{X}_n:=\frac{X_1+X_2+\cdots+X_n}{n}
$$

Determine approximately the number of resistors $n$ 
she needs to measure so that the probability that 
the sample mean differs from the true mean by more than 0.005 ohms 
is less than $1 \%$, i.e.,
$$
\mathbb{P}\left(\left|\bar{X}_n-\mu\right|>0.005\right)<0.01
$$
\begin{proof}
By linearity of expectation, 
\[
\mathbb{E}[\bar X_n]=\mu
\]
And (assmuming the \(X_i\) are independent), we have
$$
\mathrm{Var}\left(\sum_{i=1}^n X_i\right) = 
\sum_{i=1}^n \mathrm{Var}(X_i) + 
\sum_{i \neq j} \mathrm{Cov}(X_i, X_j) = n \sigma^2 = 0.25 n
$$
Thus 
\[
\mathrm{Var}(\bar X_n)=\frac{\sigma^2}{n}=\frac{0.25}{n}
\]
By Chebyshev's inequality, for any \(\varepsilon>0\),
\[
\mathbb{P}\left(|\bar X_n-\mu|>\varepsilon\right)
\le \frac{\mathrm{Var}(\bar X_n)}{\varepsilon^2}
\]
Taking \(\varepsilon=0.005\), we get
\[
\mathbb{P}\left(|\bar X_n-\mu|>0.005\right)
\le \frac{0.25/n}{(0.005)^2}
= \frac{0.25}{n\cdot 0.000025}
= \frac{10000}{n}
\]
We want this upper bound to be less than \(0.01\), 
so it is enough to require
\[
\frac{10000}{n}<0.01
\]
Therefore, she needs to measure approximately
\(
n \approx 1{,}000{,}000
\)
resistors.
\end{proof}




\section*{Problem 3}
Let $\left(X_n\right)_{n \in \mathbb{N}}$ be a sequence of random variables 
with $\mathbb{P}\left(X_n \neq 0\right)=1 / n^2$ for all $n \in \mathbb{N}$.
 Show that with probability 1, 
there exists an $n_0 \in \mathbb{N}$ such that $X_n=0$ for all $n \geq n_0$.
\begin{proof}
Let
\[
A_n := \{X_n \neq 0\}, \quad n \in \mathbb{N}
\]
Then
\(
\mathbb{P}(A_n)=\frac{1}{n^2} 
\) by assumption. 
Hence
\[
\sum_{n=1}^\infty \mathbb{P}(A_n)
= \sum_{n=1}^\infty \frac{1}{n^2}
< \infty
\]
By Borel-Cantelli lemma,
\[
\mathbb{P}\left( \limsup_{n \to \infty} A_n \right) = 0
\]
So with probability \(1\), only finitely many of the events \(A_n\) occur. In other words, with probability \(1\), there exists \(n_0 \in \mathbb{N}\) such that for all \(n \ge n_0\),
\[
A_n^c = \{X_n=0\}
\]
occurs. Equivalently,
\[
X_n=0 \qquad \text{for all } n \ge n_0
\]

Therefore, with probability \(1\), there exists \(n_0 \in \mathbb{N}\) such that \(X_n=0\) for all \(n \ge n_0\).
\end{proof}








\section*{Problem 4}
Assume that $\left(X_n\right)_{n \in \mathbb{N}}$ be a sequence of i.i.d random variables 
with density
$$
f(x)= 
\begin{cases}
\frac{1}{2 \sqrt{x}}, & \text { if } x \in(0,1), 
\\ 0, & \text { otherwise }
\end{cases}
$$
\begin{itemize}
    \item Find the distribution function of $X_1$.
    \item Let $Y_n=\min \left\{X_1, \ldots, X_n\right\}$ for any $n \in \mathbb{N}$.
     Show that $n^2 Y_n \xrightarrow{d} Y$, where $Y$ has distribution function
$$
F_Y(x)= 
\begin{cases}
0, & \text { if } x \leq 0 \\ 
1-e^{-\sqrt{x}}, & \text { otherwise }
\end{cases}
$$
\end{itemize}

\begin{proof}
\begin{itemize}
\item
For \(x \in \mathbb{R}\),
\[
F_{X_1}(x)=\mathbb{P}(X_1\le x)=\int_{-\infty}^x f(t)\,dt
=
\begin{cases}
0, & x\le 0,\\[1mm]
\displaystyle \int_0^x \frac{1}{2\sqrt{t}}\,dt=\sqrt{x}, & 0<x<1,\\[2mm]
1, & x\ge 1.
\end{cases}
\]

\item
If \(x\le 0\), then \(n^2Y_n\ge 0\), so
\[
\mathbb{P}(n^2Y_n\le x)=0
\]
Now consider \(x>0\). Then
\[
\mathbb{P}(n^2Y_n>x)=\mathbb{P}\left(Y_n>\frac{x}{n^2}\right)
\]
Since
\[
Y_n>\frac{x}{n^2}
 \iff  
X_1>\frac{x}{n^2},\dots,X_n>\frac{x}{n^2}
\]
and the \(X_i\) are independent,
\[
\mathbb{P}\left(Y_n>\frac{x}{n^2}\right)
=\left(\mathbb{P}\left(X_1>\frac{x}{n^2}\right)\right)^n
\]
For all sufficiently large \(n\), we have \(0<\frac{x}{n^2}<1\), hence
\[
\mathbb{P}\left(X_1>\frac{x}{n^2}\right)
=1-F_{X_1}\left(\frac{x}{n^2}\right)
=1-\sqrt{\frac{x}{n^2}}
=1-\frac{\sqrt{x}}{n}
\]
Therefore,
\[
\mathbb{P}(n^2Y_n>x)=\left(1-\frac{\sqrt{x}}{n}\right)^n
\]
so
\[
F_{n^2Y_n}(x)
=\mathbb{P}(n^2Y_n\le x)
=1-\left(1-\frac{\sqrt{x}}{n}\right)^n
\]
Taking \(n\to\infty\), we use the standard limit
\[
\left(1-\frac{a}{n}\right)^n \to e^{-a}
\]
with \(a=\sqrt{x}\), we obtain
\[
F_{n^2Y_n}(x)\to 1-e^{-\sqrt{x}}, \quad x>0
\]
Thus,
\[
F_{n^2Y_n}(x)\to
\begin{cases}
0, & x\le 0,\\
1-e^{-\sqrt{x}}, & x>0
\end{cases}
\]
which is exactly \(F_Y(x)\). Hence
\[
n^2Y_n \xrightarrow{d} Y
\]
\end{itemize}
\end{proof}


\section*{Problem 5}
Let $\left(X_n\right)_{n \in \mathbb{N}}$ be a sequence of random variables 
with values in $\mathbb{R}$. 
Show that there exists a sequence $\left(a_n\right)_{n \in \mathbb{N}}$ with $a_n>0$ such that
$$
\frac{X_n}{a_n} \xrightarrow{\text { a.s. }} 0
$$
For simplicity you may assume that $X_n \sim \operatorname{Exp}(1 / n)$.

Hint: For any $n \in \mathbb{N}$ construct $b_n$ such that 
$\mathbb{P}\left(\left|X_n\right| \geq b_n\right) \leq \frac{1}{2^n}$ 
and use Borel-Cantelli for the events $\left\{\left|X_n\right| / b_n \geq n\right\}$.


\begin{proof}
For each $n\in\mathbb{N}$, choose $b_n>0$ such that
$$
\mathbb{P}(|X_n|\ge b_n)\le \frac{1}{2^n}
$$
Notice such a choice is always possible, since $|X_n|$ is a well-defined 
random variable, which implies $\mathbb{P}(|X_n|\ge t)\to 0$ as $t\to\infty$.

Now define
$$a_n := n b_n > 0$$
Consider the sequence of events
$$
E_n := \left\{ \frac{|X_n|}{a_n} \ge \frac{1}{n} \right\}
$$
By the definition of $a_n$, we have
$$
E_n = \left\{ \frac{|X_n|}{n b_n} \ge \frac{1}{n} \right\} 
= \{ |X_n| \ge b_n \}
$$
It follows that
$$
\sum_{n=1}^\infty \mathbb{P}(E_n) =
 \sum_{n=1}^\infty \mathbb{P}(|X_n| \ge b_n) \le 
 \sum_{n=1}^\infty \frac{1}{2^n} < \infty
$$
By the first Borel-Cantelli lemma,
$$
\mathbb{P}(E_n \text{ infinitely often}) = 0
$$
This means that for almost all $\omega \in \Omega$, 
there exists $N(\omega) \in \mathbb{N}$ such that: for all 
$n \ge N(\omega)$, the event $E_n$ does not occur, 
i.e.,
$$\frac{|X_n(\omega)|}{a_n} < \frac{1}{n}
$$
Since $1/n \to 0$ as $n \to \infty$, it follows immediately that
$$
\frac{X_n}{a_n} \to 0 \quad \text{a.s.}
$$
If we assume $X_n \sim \mathrm{Exp}(1/n)$, 
we can provide an explicit sequence. 
Since
$$
\mathbb{P}(X_n \ge t) = e^{-t/n} \quad \text{for } t \ge 0
$$
we can choose $b_n = n^2 \log 2$. So that
$$
\mathbb{P}(X_n \ge b_n) = e^{-(n^2 \log 2)/n} = e^{-n \log 2} = \frac{1}{2^n}
$$
Then $a_n = n b_n = n^3 \log 2$, the general argument above guarantees that
$$
\frac{X_n}{a_n} = 
\frac{X_n}{n^3 \log 2} \xrightarrow{\text{a.s.}} 0
$$
This completes the proof.
\end{proof}







\section*{Problem 6}
Let $\left\{X_i\right\}_{i \geq 1}$ be i.i.d. positive integer-valued random variables 
with $0<\mathbb{E}\left[X_1\right]<\infty$. 
Interpret $X_i$ as the number of children in family $i$. 
From the first $n$ families, choose a child uniformly at random among all children. 
Let $N_n$ denote the number of children in the selected child's family. 
Show that $N_n \xrightarrow{d} X_1^*$, where $X_1^*$ has distribution

$$
\mathbb{P}\left(X_1^*=k\right)=\frac{k \mathbb{P}\left(X_1=k\right)}{\mathbb{E}\left[X_1\right]}
$$


\begin{proof}
For each \(n\in\mathbb{N}\), let
\[
S_n:=X_1+\cdots+X_n
\]
be the total number of children in the first \(n\) families.

Given \(X_1,\dots,X_n\), we choose one child uniformly 
at random among these \(S_n\) children. 
Hence, conditionally on \(X_1,\dots,X_n\), 
the probability that the chosen child comes from a family 
with exactly \(k\) children is
\[
\mathbb{P}(N_n=k \mid X_1,\dots,X_n)
=
\frac{\sum_{i=1}^n X_i \mathbf{1}_{\{X_i=k\}}}{S_n}
\]
Since \(X_i \mathbf{1}_{\{X_i=k\}} = k \mathbf{1}_{\{X_i=k\}}\), this becomes
\[
\mathbb{P}(N_n=k \mid X_1,\dots,X_n)
=
\frac{k\sum_{i=1}^n \mathbf{1}_{\{X_i=k\}}}{S_n}
=
\frac{k \cdot \frac{1}{n}\sum_{i=1}^n \mathbf{1}_{\{X_i=k\}}}{\frac{1}{n}\sum_{i=1}^n X_i}
\]

By the Strong Law of Large Numbers,
\[
\frac{1}{n}\sum_{i=1}^n \mathbf{1}_{\{X_i=k\}}
\to \mathbb{E}[\mathbf{1}_{\{X_1=k\}}]
= \mathbb{P}(X_1=k)
\quad \text{a.s.}
\]
and
\[
\frac{1}{n}\sum_{i=1}^n X_i
\to \mathbb{E}[X_1]
\quad \text{a.s.}
\]
Since \(0<\mathbb{E}[X_1]<\infty\), it follows that
\[
\mathbb{P}(N_n=k \mid X_1,\dots,X_n)
\to
\frac{k\mathbb{P}(X_1=k)}{\mathbb{E}[X_1]}
\quad \text{a.s.}   
\]

Taking expectations on both sides, and using dominated convergence because
\(
0 \le \mathbb{P}(N_n=k \mid X_1,\dots,X_n)\le 1
\),
we obtain
\[
\mathbb{P}(N_n=k)
=
\mathbb{E}\big[\mathbb{P}(N_n=k \mid X_1,\dots,X_n)\big]
\to
\frac{k\mathbb{P}(X_1=k)}{\mathbb{E}[X_1]}
\]
Thus, for every \(k\in\mathbb{N}\),
\[
\mathbb{P}(N_n=k)\to \mathbb{P}(X_1^*=k)
\]
Hence
\[
N_n \xrightarrow{d} X_1^*
\]
\end{proof}